\chapter{Incerteza e diagnósticos}\label{cap:diagnosticos}
\section{Diferenças conhecidas e incerteza auxiliar}
As diferenças de MAE nos casos e previsões observados são números conhecidos, não estimativas com erro de Monte Carlo. Para discutir sua variação em outras avaliações hipotéticas, é necessário explicitar um alvo adicional e suas hipóteses. A quantidade de componentes operacionais não os torna automaticamente uma amostra probabilística de artistas ou composições independentes.

A análise auxiliar trata componentes como unidades de dependência dentro de cada teste externo e condiciona separadamente no respectivo treino e par de modelos. Para um contraste entre configurações $a$ e $b$, define-se $D_g=\sum_{i\in g}(|e_i^{(a)}|-|e_i^{(b)}|)$ e $N_g$ como número de faixas do grupo. No fold $k$, o contraste é uma razão de totais,
\begin{equation}\label{eq:deltagrupo}
 \widehat\delta_k=\frac{\sum_{g\in k}D_g}{\sum_{g\in k}N_g}.
\end{equation}
Isso preserva o estimando por faixa; não é média simples dos MAEs por grupo. A linearização da razão conduz ao erro padrão auxiliar
\begin{equation}\label{eq:segrupo}
 \widehat{\operatorname{SE}}_k^2=
 \frac{G_k}{G_k-1}\frac{\sum_{g\in k}(D_g-\widehat\delta_kN_g)^2}
 {(\sum_{g\in k}N_g)^2}.
\end{equation}
Na sensibilidade ponderada, $D_g$ incorpora os pesos originais e $N_g$ passa a ser sua soma. A lógica admite dependência arbitrária dentro de grupos e requer independência e regularidade entre componentes em um modelo auxiliar, conforme os princípios de inferência por clusters discutidos por \citet{cameron2015}. Essas hipóteses não são demonstradas pelo desenho original de faixas e podem falhar por aliases ou relações residuais.

Foram construídos limites normais por fold e aplicada correção Bonferroni a dez quantidades, dois contrastes por cinco folds, com família separada para cada ponderação. A combinação usa médias ponderadas dos extremos por tamanho ou soma de pesos do fold. A desigualdade da união não exige independência entre folds, mas sua cobertura nominal depende da adequação dos limites marginais. Não se condiciona conjuntamente em todos os modelos: os outros treinos contêm os casos que são teste neste fold.

\begin{table}[htbp]
\centering\small
\caption{Limites auxiliares condicionais para os contrastes de interesse.}
\label{tab:incerteza}
\begin{tabular}{lrrrr}
\toprule
Contraste & Avaliação & Delta & Inferior & Superior \\
\midrule
C - B & Iguais & -0,23736 & -0,31400 & -0,16073 \\
C - P & Iguais & -0,01340 & -0,03027 & 0,00347 \\
C - B & Originais & -0,22861 & -0,30454 & -0,15268 \\
C - P & Originais & -0,01234 & -0,02919 & 0,00450 \\
\bottomrule
\end{tabular}
\fonte{Elaboração própria a partir dos agregados auditados do estudo.}
\nota{C: Combinado; B: Baseline; P: Profundo. 95\% nominais sob modelo auxiliar de componentes, com Bonferroni por família de dez marginais. Não são ICs da pipeline integral.}
\end{table}

\figura{limites_condicionais}{Contrastes e limites auxiliares condicionais.}{fig:limites}{Análise exploratória por componentes/folds, com 95\% nominais e Bonferroni; C significa Combinado. Não representa a incerteza integral de seleção e treinamento.}
\FloatBarrier

Os limites da Tabela~\ref{tab:incerteza} mantêm direção favorável para Combinado--Baseline, mas incluem zero para Combinado--Profundo. O segundo ganho é inconclusivo mesmo nesse cenário auxiliar. O alvo combinado é uma média de riscos condicionais dos cinco pares ajustados em avaliações hipotéticas compatíveis, e não o risco de todo o algoritmo refeito em novas amostras ou o do modelo ajustado em todas as faixas. As escolhas históricas de hipóteses e famílias não são apagadas por Bonferroni; a análise continua exploratória.

Métodos assintóticos de inferência para CV podem ser defensáveis sob condições de estabilidade \citep{bayle2020}. Porém, estabilidade individual dos modelos não basta para a estabilidade relativa de uma comparação \citep{bayle2026}. Não se presume que o nosso Ridge viole esses resultados, mas tampouco que a simples escolha dessa família valide seus intervalos. \citet{bates2024} estudam um desenho aninhado para estimar variabilidade da CV; ele não é automaticamente equivalente ao aninhamento 5$\times$3 aqui usado para tuning.

Um bootstrap pareado de componentes com previsões congeladas avaliaria composição de casos, não toda a incerteza de treinamento. Para incorporar seleção e reajuste seria necessário outro desenho, com pré-processamento, tuning e modelos refeitos em cada réplica e manutenção das cópias de um componente na mesma partição. Nenhum refit desse tipo foi executado ou é pressuposto nos limites apresentados.

\section{Influência de componentes e subgrupos}
A retirada de um componente apenas da avaliação preservou o sinal negativo dos cinco contrastes, tanto sem pesos quanto com pesos originais. A Tabela~\ref{tab:remocao} mostra os extremos condicionais. Esse resultado indica que um único componente não é suficiente para inverter a direção dos contrastes observados com os ajustes fixos. Não é validação leave-one-group-out com refit e seus extremos não são limites de confiança.

\begin{table}[htbp]
\centering\small
\caption{Influência da retirada de um componente apenas da avaliação.}
\label{tab:remocao}
\begin{tabular}{lrrr}
\toprule
Contraste & Delta & Menor delta & Maior delta \\
\midrule
Tradicional - Baseline & -0,05679 & -0,05798 & -0,05570 \\
Profundo - Baseline & -0,22396 & -0,22456 & -0,22254 \\
Combinado - Baseline & -0,23736 & -0,23806 & -0,23588 \\
Combinado - Tradicional & -0,18057 & -0,18141 & -0,17916 \\
Combinado - Profundo & -0,01340 & -0,01381 & -0,01316 \\
\bottomrule
\end{tabular}
\fonte{Elaboração própria a partir dos agregados auditados do estudo.}
\nota{Modelos e previsões fixos. Extremos não são limites de confiança nem validação com refit.}
\end{table}


Os diagnósticos por subgrupo publicaram métricas apenas em células com pelo menos 100 faixas e 20 componentes; outras tiveram contagens e motivo de omissão, sem excluir casos da avaliação global. Foram 116 células com métricas e 77 com métricas omitidas. Esses limiares controlam a apresentação de células pequenas e não garantem precisão inferencial.

\begin{table}[htbp]
\centering\small
\caption{Desempenho por status do idioma estimado.}
\label{tab:idiomas}
\begin{tabular}{lrrrr}
\toprule
Status & n & Baseline & Profundo & Combinado \\
\midrule
Inglês estimado & 56.908 & 7,6840 & 7,5418 & 7,5278 \\
Não inglês estimado & 20.447 & 10,2036 & 9,7491 & 9,7343 \\
Idioma incerto & 180 & 9,4360 & 9,5398 & 9,8564 \\
\bottomrule
\end{tabular}
\fonte{Elaboração própria a partir dos agregados auditados do estudo.}
\nota{Idioma automático não é rótulo humano. Diagnóstico exploratório sem excluir casos.}
\end{table}

A Tabela~\ref{tab:idiomas} mostra que o ganho não é uniforme. No status incerto, com 180 casos, o MAE do Combinado foi maior que o do Baseline. Entre os 76 gêneros com métricas publicadas, 59 favoreceram o Combinado e 17 não. Os resultados são exploratórios, com sobreposição de componentes entre categorias; não são uma bateria confirmatória de testes. O idioma automático também pode errar e não justifica redefinir a população após a comparação.

\figura{erro_idioma}{MAE por status de idioma estimado.}{fig:idiomas}{Diagnóstico agregado das mesmas previsões externas, sem alterar o domínio analítico ou o detector de idioma.}
\FloatBarrier

\section{Cauda de popularidade e distribuição dos resíduos}
A dificuldade na cauda superior é expressiva. Os 128 casos com $y\ge80$ pertencem a 96 componentes e apresentam valores de 80 a 96; 116 estão em 80--84, oito em 85--89 e quatro em 90--100. Não se trata de 128 observações próximas de cem. Todas foram subestimadas pelas quatro configurações, conforme a Tabela~\ref{tab:cauda}.

\figura{erro_popularidade}{Erro por banda de popularidade observada no Baseline e no Combinado.}{fig:bandaserro}{Diagnósticos da comparação recente; contagens indicadas na figura. As bandas são definidas pelo alvo, e não pelo valor previsto.}
\FloatBarrier

\begin{table}[htbp]
\centering\small
\caption{Previsões nas 128 faixas com popularidade observada de 80 a 100.}
\label{tab:cauda}
\begin{tabular}{lrrrr}
\toprule
Configuração & Previsão média & MAE & RMSE & Subestimadas \\
\midrule
Baseline & 50,8831 & 30,9294 & 32,7434 & 128 \\
Tradicional & 51,3634 & 30,4491 & 32,2431 & 128 \\
Profundo & 52,1361 & 29,6764 & 31,6598 & 128 \\
Combinado & 52,4414 & 29,3711 & 31,3288 & 128 \\
\bottomrule
\end{tabular}
\fonte{Elaboração própria a partir dos agregados auditados do estudo.}
\nota{Média observada: 81,8125. Há 96 componentes; a faixa efetivamente observada vai de 80 a 96.}
\end{table}

No Combinado, a previsão média foi 52,4414 para popularidade média de 81,8125, com previsões entre 24,9681 e 77,9977. Definindo agora o resíduo diagnóstico como $r_i=\widehat y_i-y_i$, valor negativo significa subestimação. Sua mediana na cauda foi $-28{,}7755$, P25 $-37{,}2444$ e P75 $-22{,}4381$; o desvio populacional descritivo foi 10,9011. Houve 101 subestimações de pelo menos 20 pontos e 22 de pelo menos 40, de modo que o MAE elevado não é produzido apenas por um ou dois casos isolados.

\figura{residuos_cauda}{Contagens de resíduos do Combinado na cauda observada ($n=128$).}{fig:residuoscauda}{Previsões existentes; resíduos iguais à previsão menos o observado. Barras são contagens por intervalo, não densidades.}
\FloatBarrier

A cauda representa 0,1651\% das faixas, 0,5975\% do erro absoluto total e 1,4591\% do erro quadrático total do Combinado. Sua contribuição aditiva ao MAE global, mantendo o denominador de 77.535, é 0,04849 ponto. A melhora frente ao Baseline nessa cauda corresponde a apenas 1,0838\% da melhora total de MAE. O bom desempenho médio relativo convive, portanto, com dificuldade de reproduzir os valores mais altos.

A compressão das previsões é compatível com regularização e com informação insuficiente para distinguir a cauda, mas esses diagnósticos não identificam causalmente o mecanismo. Condicionar no próprio alvo alto pode produzir resíduos negativos mesmo sem erro de código. Isso não é a mesma análise que calibração condicional na previsão, $E(Y\mid\widehat Y)$. No domínio inteiro, o Combinado não emitiu previsão de pelo menos 80.

\section{Previsões fora do suporte e sensibilidade ao clipping}
No Combinado foram observadas 615 previsões negativas, em 317 componentes, com mínimo $-32{,}7066$ e mediana $-1{,}7873$. Nesses casos, a resposta tem média 5,3122 e mediana 3; 133 são zero e apenas um caso tem popularidade pelo menos 40. A extrapolação ocorre principalmente na borda inferior. Não é evidência, por si só, de corrupção de embeddings.

Como sensibilidade separada, calculou-se $\widetilde y=\min(100,\max(0,\widehat y))$, sem substituir as previsões principais. Se $y\in[0,100]$, a projeção não aumenta o erro absoluto ou quadrático ponto a ponto. Por isso, uma melhora após clipping é propriedade da restrição de suporte, não comprovação de nova informação textual.

\begin{table}[htbp]
\centering\small
\caption{Projeção de suporte como sensibilidade pós-resultado.}
\label{tab:clipping}
\begin{tabular}{lrrrr}
\toprule
Configuração & Abaixo de 0 & Acima de 100 & MAE original & MAE projetado \\
\midrule
Baseline & 537 & 0 & 8,35248 & 8,33781 \\
Tradicional & 584 & 1 & 8,29569 & 8,28013 \\
Profundo & 610 & 0 & 8,12852 & 8,11055 \\
Combinado & 615 & 0 & 8,11512 & 8,09701 \\
\bottomrule
\end{tabular}
\fonte{Elaboração própria a partir dos agregados auditados do estudo.}
\nota{Projeção fixa em [0,100], sem substituir as previsões ou métricas principais.}
\end{table}

A ordenação dos modelos permanece a mesma. O MAE do Combinado passa a 8,0970 nessa sensibilidade, redução de 0,01811 ponto. A projeção não modifica nenhuma das 128 previsões da cauda, pois já estavam em [0,100]. Trata-se de análise posterior aos resultados, que não altera o estimador declarado nem a comparação principal.
